\subsection{STRUCTURE OF FILTERED BIGEBRAS IN CHARACTERISTIC 0}

In this no. we continue to assume that K is a field of characteristic 0.

If E is a bigebra, the canonical injection \(P(E) \rightarrow E\) can be extended to a bigebra morphism \(f_{E}: U(P(E)) \rightarrow E\) (no. 4, Proposition 3).

THEOREM 1. Let E be a cocommutative bigebra.

\begin{enumerate}
\def\labelenumi{(\alph{enumi})}
\item
  The bigebra morphism \(f_{\mathbf{E}} \colon \mathrm{U}(\mathrm{P}(\mathbf{E})) \to \mathbf{E}\) is injective.
\item
  If there exists on E a filtration compatible with its bigebra structure (no. 3, Definition 2), the morphism \(f_{\mathrm{E}}\) is an isomorphism.
\end{enumerate}

(In case (b), the bigebra E is therefore identified with the enveloping bigebra of the Lie algebra of its primitive elements.)

Let \(c_{\mathbb{E}}\) (resp. \(\varepsilon_{\mathbb{E}}\) ) be the coproduct (resp. counit) of E. We write \(\mathfrak{g} = \mathrm{P}(\mathrm{E})\) ; let \((e_i)_{i \in \mathrm{I}}\) be a basis of the vector K-space \(\mathfrak{g}\) , where the indexing set I is totally ordered, and let \((e_\alpha)_{\alpha \in \mathbf{N}^{(I)}}\) be the basis introduced in the preceding no. We write \(\mathbf{X}_\alpha = f_{\mathbb{E}}(e_\alpha)\) for \(\alpha \in \mathbf{N}^{(I)}\) . By (15) and (16), we have:

\begin{enumerate}
\def\labelenumi{(\arabic{enumi})}
\setcounter{enumi}{16}
\tightlist
\item
\end{enumerate}

\[
\varepsilon_ {E} (X _ {0}) = 1, \quad \varepsilon_ {E} (X _ {\alpha}) = 0 \quad \text { for } | \alpha | \geqslant 1,\tag{18}
\]

\[
c _ {E} (X _ {\alpha}) = \sum_ {\beta + \gamma = \alpha} X _ {\beta} \otimes X _ {\gamma} \quad \text { for } \alpha \in \mathbf {N} ^ {(I)},
\]

since \(f_{E}\) is a cogebra morphism.

We show that \(f_{E}\) is injective. This results from the following lemma:

Lemma 2. Let V be a vector space, E a cogebra and \(f\colon\mathsf{S}(\mathrm{V})\to\mathrm{E}\) a cogebra morphism. If the restriction of \(f\) to \(\mathsf{S}^{\circ}(\mathrm{V})+\mathsf{S}^{1}(\mathrm{V})\) is injective, then \(f\) is injective.

Let \(n \geqslant 0\) ; we write \(\mathbb{S}_n = \sum_{i \geqslant n} \mathbb{S}^i(\mathbf{V})\) and \(c_s\) for the coproduct of \(\mathbb{S}(\mathbf{V})\) and show by induction on \(n\) that \(f \mid \mathbb{S}_n\) is injective. Since the assertion is trivial for \(n = 0\) and \(n = 1\) , we assume that \(n \geqslant 2\) and let \(u \in \mathbb{S}_n\) be such that \(f(u) = 0\) . Then

\[
\begin{array}{r l} {0 = c _ {\mathbf {E}} (f (u))} & {= (f \otimes f) (c _ {\mathbf {s}} (u))} \\ & {= f (u) \otimes 1 + 1 \otimes f (u) + (f \otimes f) (c _ {\mathbf {s}} ^ {+} (u))} \\ & {= (f \otimes f) (c _ {\mathbf {s}} ^ {+} (u)).} \end{array}
\]

As \(c_{\mathbb{S}}^{+}(u) \in \mathbb{S}_{n-1} \otimes \mathbb{S}_{n-1}\) , by (11) the induction hypothesis shows that \(u\) is a primitive element of \(\mathbb{S}(\mathrm{V})\) , hence is of degree 1 (no. 5, Remark 4) and hence is zero, since \(f \mid \mathbb{S}^1(\mathrm{V})\) is injective.

It follows in particular that the family \((X_{\alpha})\) is free.

We show that \(f_{\mathbb{E}}\) is surjective if \(\mathrm{E}\) has a filtration compatible with its bigebra structure. Let \((\mathrm{E}_n)_{n\geqslant 0}\) be such a filtration and write \(\mathrm{E}_n^+ = \mathrm{E}_n\cap \mathrm{Ker}(\epsilon_{\mathbb{E}})\) . We show by induction on \(n\) that \(\mathrm{E}_n^+\) is contained in the image of \(f_{\mathbb{E}}\) . As \(\mathrm{E} = \mathrm{K}.1 + \bigcup_{n > 0}\mathrm{E}_n^+\) , this will imply the surjectivity of \(f_{\mathbb{E}}\) . The assertion is trivial for \(n = 0\) and follows from the Corollary to Proposition 6 of no. 3 for \(n = 1\) ; suppose henceforth that \(n\geqslant 2\) and let \(x\in \mathrm{E}_n^+\) . By Proposition 6 of no. 3,

\[
c _ {\mathrm{E}} ^ {+} (x) \in \sum_ {i = 1} ^ {n - 1} \mathrm{E} _ {i} ^ {+} \otimes \mathrm{E} _ {n - i} ^ {+}
\]

and by the induction hypothesis there exist scalars \(\lambda_{\alpha, \beta}\) , where \(\alpha, \beta\) are in \(\mathbf{N}^{(I)}\) , which are zero except for a finite number, such that

\[
c _ {\mathbb {E}} ^ {+} (x) = \sum_ {\alpha , \beta \neq 0} \lambda_ {\alpha , \beta} \mathrm{X} _ {\alpha} \otimes \mathrm{X} _ {\beta}.\tag{19}
\]

Hence by formula (18)

\[
(c _ {\mathrm{E}} ^ {+} \otimes \operatorname{Id} _ {\mathrm{E}}) (c _ {\mathrm{E}} ^ {+} (x)) = \sum_ {\alpha , \beta , \gamma \neq 0} \lambda_ {\alpha + \beta , \gamma} \mathbf {X} _ {\alpha} \otimes \mathbf {X} _ {\beta} \otimes \mathbf {X} _ {\gamma}
\]

\[
(\mathrm{Id} _ {\mathtt {E}} \otimes c _ {\mathtt {E}} ^ {+}) (c _ {\mathtt {E}} ^ {+} (x)) = \sum_ {\alpha , \beta , \gamma \neq 0} \lambda_ {\alpha , \beta + \gamma} \mathbf {X} _ {\alpha} \otimes \mathbf {X} _ {\beta} \otimes \mathbf {X} _ {\gamma}.
\]

By Proposition 3 of no. 1 and the linear independence of the \(X_{\alpha}\) it follows that

\[
\lambda_ {\alpha + \beta , \gamma} = \lambda_ {\alpha , \beta + \gamma} \quad \text { for } \alpha , \beta , \gamma \text { in } \mathbf {N} ^ {(I)} - \{0 \}.\tag{20}
\]

On the other hand, the coproduct \(c_{E}\) is cocommutative; the same argument as above implies

\[
\lambda_ {\alpha , \beta} = \lambda_ {\beta , \alpha} \quad \text { for } \alpha , \beta \text { in } \mathbf {N} ^ {(I)} - \{0 \}.\tag{21}
\]

Suppose that there exists a family of scalars \((\mu_{\alpha})\) with \(|\alpha| \geqslant 2\) , such that

\[
\mu_ {\alpha + \beta} = \lambda_ {\alpha , \beta} \quad \text { for } \alpha , \beta \text { in } \mathbf {N} ^ {(I)} - \{0 \}.\tag{22}
\]

Then

\[
c _ {\mathrm{E}} ^ {+} (x) = \sum_ {\alpha , \beta \neq 0} \mu_ {\alpha + \beta} \mathrm{X} _ {\alpha} \otimes \mathrm{X} _ {\beta} = \sum_ {| \gamma | \geqslant 2} \mu_ {\gamma} c _ {\mathrm{E}} ^ {+} (\mathrm{X} _ {\gamma})
\]

by formula (18), hence \(y = x - \sum_{|\gamma| \geq 2} \mu_\gamma X_\gamma\) is primitive and hence belongs to \(P(E) \subset \operatorname{Im}(f_E)\) . Hence

\[
x = y + \sum_ {| \gamma | \geqslant 2} \mu_ {\gamma} f _ {\mathrm{E}} (e _ {\gamma}) \in \operatorname{Im} (f _ {\mathrm{E}}).
\]

The proof will therefore be complete when we have proved the following lemma:

Lemma 3. If a family of scalars \((\lambda_{\alpha, \beta})\) of finite support (with \(\alpha, \beta\) in \(\mathbf{N}^{(I)} - \{0\}\) ) satisfies relations (20) and (21), there exists a family \((\mu_{\alpha})_{|\alpha| \geqslant 2}\) of finite support such that \(\mu_{\alpha + \beta} = \lambda_{\alpha, \beta}\) for \(\alpha, \beta\) non-zero.

It suffices to prove that

\[
\alpha + \beta = \gamma + \delta\tag{23}
\]

implies \(\lambda_{\alpha, \beta} = \lambda_{\gamma, \delta}\) for \(\alpha, \beta, \gamma, \delta\) non-zero. By Riesz's Decomposition Lemma (Algebra, Chapter VI, § 1, no. 10, Theorem 1) there exist \(\pi, \rho, \sigma, \tau\) in \(\mathbf{N}^{(I)}\) such that

\[
\alpha = \pi + \sigma , \quad \beta = \rho + \tau , \quad \gamma = \pi + \rho , \quad \delta = \sigma + \tau .
\]

Suppose \(\pi \neq 0\) ; as \(\sigma + \beta = \rho + \delta\) , relation (20) implies

\[
\lambda_ {\alpha , \beta} = \lambda_ {\pi + \sigma , \beta} = \lambda_ {\pi , \sigma + \delta} = \lambda_ {\pi , \rho + \delta} = \lambda_ {\pi + \rho , \delta} = \lambda_ {\gamma , \delta}.
\]

If on the other hand \(\pi = 0\) , then \(\beta = \gamma + \tau\) and \(\delta = \alpha + \tau\) , whence

\[
\lambda_ {\alpha , \beta} = \lambda_ {\alpha , \gamma + \tau} = \lambda_ {\alpha + \tau , \gamma} = \lambda_ {\delta , \gamma}
\]

by (20), but also \(\lambda_{\delta, \gamma} = \lambda_{\gamma, \delta}\) by (21), whence \(\lambda_{\alpha, \beta} = \lambda_{\gamma, \delta}\) .

